\chapter{Biomass}

\section{Introduction}

Biomass — organic matter derived from living or recently living organisms — represents the oldest energy source used by humanity and remains today a central, if contested, pillar of the global energy system. This document synthesises the key technical, statistical, environmental, and philosophical dimensions of biomass as presented in a lecture on biomass and the energy transition, with particular attention to its sustainability. The core argument running through the material is that biomass itself is neither inherently sustainable nor unsustainable: what determines its environmental and social value is the system of management and exploitation within which it is embedded.

\section{What Is Biomass?}

Biomass designates any organic material of biological origin used for energy purposes. It encompasses a wide and heterogeneous range of feedstocks:

\begin{itemize}
    \item \textbf{Forestry biomass}: stemwood, logging residues (treetops, branches), and material from thinning operations.
    \item \textbf{Wood industry by-products}: sawdust, bark, and chips generated during the processing of timber into lumber or panels.
    \item \textbf{Energy crops}: lignocellulosic crops (e.g., miscanthus, short-rotation coppice), oil crops (sunflower, palm, rapeseed), and sugar/starch crops (corn, wheat, sugar cane, sugar beet) used for bioethanol or biodiesel production.
    \item \textbf{Agricultural residues}: straw and other crop residues left after harvest.
    \item \textbf{Animal and agro-industrial residues}: manure, slaughterhouse by-products, and organic food processing waste.
    \item \textbf{Waste streams}: the organic fraction of municipal solid waste (40–60\% of the total), demolition wood, paper and cardboard, and sewage sludge.
\end{itemize}

These feedstocks can be converted through multiple routes — combustion, transesterification, hydrolysis and fermentation, gasification, pyrolysis, and anaerobic digestion — into a variety of end-uses: heat, electricity, liquid transport fuels, and chemical feedstocks.

\section{The Current Role of Biomass in Energy Systems}

\subsection{Global and European Context}

Historically, traditional biomass was the dominant global energy source until the industrial revolution, after which fossil fuels rapidly took over. Today, modern bioenergy has regained importance. In the European Union, biomass is the \textbf{leading renewable energy source}, accounting for the largest share of primary renewable energy production. Nevertheless, the overall consumption mix in Europe remains heavily dominated by fossil fuels, with biomass representing a comparatively modest share of total primary energy consumption.

A similar structure is observed in Belgium, where the consumption pattern mirrors that of the EU27: fossil fuels dominate, with biomass playing a significant but secondary role among renewables.

\subsection{End-Uses of Biomass in Europe}

The dominant end-use of biomass in Europe is \textbf{heat production}, particularly for residential and industrial applications. Biofuels for transport represent approximately 13\% of biomass use. Bioelectricity production, while growing, is a smaller share. Among the different biomass types used in the EU27, solid biomass is overwhelmingly dominant, followed at a much lower level by liquid biofuels, biogas, and renewable municipal waste.

\subsection{Feedstock Composition}

In Europe in 2020, \textbf{forestry products} represent by far the largest feedstock category for bioenergy (approximately 1,200 TWh), with agricultural products and other waste each contributing roughly 250 TWh. Within solid biomass, pellets represent only about 10\% of the total; the remaining 90\% consists of other solid biomass forms such as firewood and chips. Within biofuels, biodiesels account for 76\%, bioethanol for 17\%, and other biofuels for 7\%.

According to a JRC (2021) analysis, the origin of woody biomass used for bioenergy in the EU (2015 data) breaks down as follows: 49\% is secondary woody biomass (largely domestic solid by-products at 21\%, and black liquor at 15\%), while 37\% is primary wood. Of this primary wood, roughly 47\% is stemwood and 53\% consists of other forest components such as treetops and branches. Additionally, 24\% is fuelwood and 4\% industrial roundwood. This means that a substantial part of wood used for energy is a co-product of forestry management operations, including thinnings, rather than wood harvested primarily for energy.

In summary, wood-energy in the EU is a mix of approximately 50\% wood extracted directly from forests (largely as a co-product of timber forestry management) and 50\% by-products from the wood processing industry. Sawing logs typically yields only 45–60\% of usable timber volume, leaving a large stream of residues available for energy use.

\section{Sustainability of Biomass}

\subsection{The Debate on Liquid Biofuels}

Liquid biofuels — particularly first-generation biofuels based on food and feed crops — have been at the centre of significant controversy. The Belgian case of Biowanze illustrates the complexity: this facility processes around 800,000 tonnes of wheat annually (50\% from Belgium, 30\% from France, 10\% from Germany, 10\% from other EU countries) to produce bioethanol alongside food-grade co-products (proteins, liquid CO\textsubscript{2}) and green energy.

Life cycle assessment (LCA) of Belgian wheat bioethanol shows that compared to gasoline (per MJ of energy delivered), bioethanol performs better on \textbf{climate change} (approximately 0.04 kg CO\textsubscript{2}eq vs. 0.084) and on fossil fuel depletion. However, it performs worse on \textbf{particulate matter formation}, \textbf{acidification}, \textbf{terrestrial and freshwater eutrophication}, \textbf{land use}, and \textbf{mineral depletion}. The verdict therefore depends critically on which environmental impact categories are prioritised.

The Energy Return on Energy Invested (EROI) — defined as the ratio of energy output to energy input across the life cycle — provides a complementary physical lens. The EROI of first-generation bioethanol is particularly low, often close to 1. For wheat-based bioethanol, reported values range from 0.69 to 2.31 depending on system boundaries, co-product allocation methods, and regional conditions. This places corn- and wheat-based ethanol near the bottom of the so-called "net energy cliff," far below fossil fuels (EROI $\sim$20–100 historically) and below other renewables such as wind and photovoltaic energy.

EU policy has responded to these concerns: the revised Renewable Energy Directive established decreasing caps for first-generation biofuels (from 7\% in 2021 to 3.8\% in 2030), specifically to limit emissions linked to indirect land-use change.

\subsection{The Carbon Cycle of Forests}

A critical issue for the sustainability of woody biomass is the accounting of carbon stocks. In EU27 forests (2020 data), carbon is distributed as follows: 37\% in above-ground biomass, 10\% in below-ground biomass, 2\% in deadwood, 7\% in litter, and \textbf{44\% in soil}. This means that the majority of forest carbon is not in the trees themselves but in the soil, a stock that is particularly sensitive to management practices and residue extraction.

At the scale of an individual tree stand, harvesting inevitably reduces the carbon stock. However, a forest is not managed at the stand scale. At the landscape scale, cutting trees in one location while replanting and managing other areas over the long term can stabilise or even increase total carbon stocks. The sustainability of carbon management therefore depends fundamentally on how the forest is managed — rotation length, residue extraction rates, species composition, and regeneration practices — rather than on the simple fact of harvesting.

\subsection{The Management and Exploitation System}

The lecture's central thesis is encapsulated in the following framing: it is not biomass itself that is or is not sustainable — it is the \textbf{system of management and exploitation}. This invites four key questions:

\begin{enumerate}
    \item What system of forestry or agricultural management is currently in place?
    \item What are the risks to sustainability under current logics of productivity and profit?
    \item What alternatives exist, and how can coherent changes be implemented?
    \item Are we even asking the right questions?
\end{enumerate}

\section{Biomass Potential and Its Variability}

Estimates of the sustainable bioenergy potential of Europe vary enormously across studies, depending on the type of potential considered (theoretical, technical, economic, ecologically sustainable, or implementable) and on key assumptions:

\begin{itemize}
    \item Dietary choices and associated land availability for energy crops;
    \item Agricultural landscape structure and scale of biogas facilities;
    \item Degree of mechanisation and use of inputs (fertilisers, irrigation);
    \item Competing uses for biomass (food, feed, materials, biodiversity conservation) and sustainable extraction rates.
\end{itemize}

Reviewed European studies yield central estimates ranging roughly from 2,100 to 5,200 TWh/year depending on the scenario, compared to a 2019 actual production of approximately 1,540 TWh. This variability is not merely technical: each estimate embeds a specific vision of the agricultural, forestry, and waste management systems, with attendant social, economic, and environmental implications. Choosing one estimate over another in an energy system model is therefore a \textbf{political act}, not a neutral technical choice, as it shapes the resulting modelled system and the policy recommendations derived from it.

Energy system models such as EnergyScope Pathway — which integrates 28 resources, 113 conversion technologies, and 12 demand categories — show that meeting a fossil-and-nuclear-free European energy system by 2050 would require a very large scale-up of biomass alongside wind and solar. In such scenarios, woody biomass could represent around 16\% of gross available energy, and wet biomass around 7\%.

\section{Broader Dimensions: Technology, Society, and the Energy Transition}

\subsection{The Limits of a Purely Technical Approach}

A purely engineering perspective, however sophisticated, risks overlooking essential dimensions of the energy transition. As the statistician George Box noted, all models are wrong but some are useful; the key is transparency about assumptions and limitations. Energy system models reduce complex phenomena to equations, and the selection of variables and boundaries can introduce biases and mask important interactions — social, ecological, and cultural.

This limitation becomes critical when considering the scale of change required. European primary energy consumption has remained broadly stable over the past two decades, dominated by fossil fuels. Transitioning to a renewable system within 30 years appears practically unfeasible without a genuine reduction in energy demand, which in turn requires a change in mentality and not just a change in technology.

\subsection{Technology Is Not Neutral}

Drawing on thinkers including Lewis Mumford, François Jarrige, and Bernard Charbonneau, the lecture argues that technical choices are never socially or politically neutral. Engineers, historically servants of military and industrial power, have been oriented primarily toward the transition of tools. But as the philosopher Ivan Illich argued (cited via Geerts 2017), any energy transition is simultaneously a transition of tools \textit{and} a transition of thinking.

\subsection{Biomass and the Human–Nature Dualism}

The debate around bioenergy is particularly intense, in part, because it surfaces deep cultural tensions about humanity's relationship with nature. Burning forests for energy makes visible — and locally tangible — the consequences of energy consumption, confronting communities with a sight many prefer to leave unseen. This visibility has, paradoxically, contributed to raising awareness of sustainable forestry practices and environmentalism.

The lecture draws on the philosopher Dominique Bourg's concept of \textit{reflective monism}: the idea that humans and non-humans belong to the same living world, with the human specificity being the capacity for reflection. Under this framing, treating biomass purely as an instrumental resource for human productivity — without attending to the complex ecological interactions on which living systems depend — is antithetical to sustainability. Imposing purely productivist logics obscures the networks of relationships essential to ecosystem integrity.

Biomass thus becomes a \textit{heuristic concept}: a lens through which to revisit the forgotten constraint of renewability, to redefine our ethical sensibility toward the non-human, and to open necessary democratic debates about what constitutes useful versus superfluous energy services. As philosopher Jonathan Parsons writes: every energy system is first and finally a social system, and a transition from one to another implies a sociopolitical transition as much as a technical one.

\section{Take-Away Messages}

The following conclusions synthesise the key points of the lecture:

\begin{itemize}
    \item Biomass can be a genuine ally in the energy transition, but sustainable implementation is essential to avoid systems driven purely by profit and productivity.
    \item It is a key energy source — particularly for heat in the EU — but one subject to limits and constraints that must be respected.
    \item The sustainable potential of biomass is highly variable and depends on interdisciplinary parameters (ecological, social, economic, dietary); estimating it is inherently a political exercise.
    \item The energy transition is fundamentally a political project that requires interdisciplinary analysis and broad democratic deliberation, not only technical optimisation.
    \item A genuine transition requires not only a change of energy sources but a change in our conceptions of nature, energy, and the relationship between human societies and ecosystems.
\end{itemize}



\section{Podcast -- Professeur Julien Blondeau, VUB}

\subsubsection{Comment et pourquoi change-t-on une centrale charbon en centrale biomasse ?}

La biomasse partage de nombreuses caractéristiques avec le charbon : c'est un combustible solide, stockable, aux bonnes propriétés de combustion. La conversion est donc techniquement envisageable. Cependant, plusieurs différences compliquent la transition.

Premièrement, la densité énergétique de la biomasse est inférieure à celle du charbon~: il faut donc en brûler davantage pour produire la même quantité d'électricité, ce qui peut limiter la puissance de la centrale.

Deuxièmement, la biomasse contient plus de matières volatiles que le charbon, ce qui modifie le comportement lors de la combustion (davantage de flammes, moins de résidu charbonneux).

Troisièmement, la taille des particules est plus grande~: broyer de la biomasse est plus difficile que broyer du charbon, et les particules obtenues (1--2~mm) sont dites \textit{thermiquement épaisses} selon le nombre de Biot. Cela signifie qu'on ne peut pas négliger les phénomènes internes à la particule lors de la modélisation de la combustion, contrairement aux particules de charbon (inférieures à 0,1~mm).

Enfin, les cendres de biomasse ont une température de fusion plus basse que celles du charbon, ce qui peut entraîner leur dépôt et l'encrassement des parois de la chaudière.

En Belgique, la centrale d'Awirs (près de Liège), opérée par Engie, a été la première au monde à être convertie à 100~\% à la biomasse, en 2005. Elle a fermé en 2020. La centrale de Rodenhuize (près de Gand, environ 200~MW électriques) était un autre exemple, avant sa fermeture suite à une décision du gouvernement régional flamand de ne plus lui accorder de subsides.

\subsubsection{Quelles sont les émissions problématiques de la biomasse ?}

La combustion de biomasse engendre, outre le CO\textsubscript{2}, d'autres types d'émissions~:
\begin{itemize}
    \item des \textbf{poussières} (cendres fines) ;
    \item des \textbf{oxydes d'azote} (NO\textsubscript{x}).
\end{itemize}

La biomasse contient très peu de soufre, ce qui l'exempte des émissions d'oxydes de soufre caractéristiques du charbon.

Dans les grandes centrales, ces polluants sont captés par des systèmes de filtration performants, déjà en place pour le charbon et largement suffisants pour la biomasse. En revanche, à l'échelle domestique (poêles, chaudières à granulés), l'absence de tels filtres rend les émissions de particules plus préoccupantes. La recherche doit donc progresser pour améliorer les performances des petites installations en termes de qualité de combustion et de réduction des émissions.

\subsubsection{Comment définir la biomasse ?}

La biomasse désigne toute source d'énergie ou tout combustible issu d'un organisme \textit{récemment} vivant. Le qualificatif « récemment » est essentiel pour exclure le charbon, qui est en réalité de la biomasse fossilisée sur des millions d'années.

Cette définition englobe un large spectre de ressources~:
\begin{itemize}
    \item le \textbf{bois} (mort ou résidus de l'industrie forestière) ;
    \item les \textbf{produits et résidus agricoles} ;
    \item la fraction \textbf{biodégradable des déchets ménagers} (environ 50~\% du total) ;
    \item les \textbf{boues de stations d'épuration} ;
    \item les \textbf{déchets animaux} ;
    \item les \textbf{résidus de l'industrie agroalimentaire}.
\end{itemize}

Aujourd'hui, le bois représente environ 70~\% de la biomasse utilisée en Europe, en raison de sa facilité d'utilisation dans de nombreuses technologies de conversion, notamment la combustion directe.

\subsubsection{Toutes les biomasses ne sont pas équivalentes}

Les différentes sources de biomasse ne sont pas interchangeables~: leur pertinence dépend de la technologie de conversion envisagée. Certaines biomasses se prêtent mieux à certains procédés.

Les déchets agricoles, par exemple, sont plus difficiles à valoriser~: leur humidité élevée complique la combustion et renchérit le transport (transporter de l'eau est coûteux en énergie et en argent). Un séchage préalable est envisageable mais lui-même énergivore. Par ailleurs, leurs cendres posent davantage de problèmes de fusion dans les chaudières.

À mesure que des usages à plus haute valeur ajoutée seront développés pour les biomasses de qualité (bioplastiques, biomatériaux, bois d'œuvre), le secteur de la bioénergie devra se tourner vers des sources plus difficiles à convertir, nécessitant des technologies plus robustes ou des étapes de prétraitement adaptées. Ces \textit{routes technologiques} sont un axe de recherche et d'optimisation important.

\subsubsection{D'où provient le bois-énergie utilisé en Belgique ?}

Dans les grandes centrales belges, le bois utilisé ne provient pas des forêts ardennaises, mais principalement de l'étranger~: États-Unis, Canada, pays scandinaves, Russie. Ces régions disposent de vastes forêts et d'une industrie forestière génératrice de résidus importants.

Il est essentiel de préciser qu'il ne s'agit pas de déforestation~: on valorise les \textbf{résidus de l'industrie forestière durable} (branches, sciures, chutes) qui, historiquement, étaient brûlés sur place en bord de route faute d'autre débouché.

Concernant l'empreinte carbone du transport, le professeur Blondeau souligne que les calculs sont réglementairement encadrés par la Commission européenne~: les émissions de CO\textsubscript{2} liées au transport ne peuvent dépasser un certain seuil. En pratique, le transport transocéanique par supertanker (par exemple depuis les côtes canadiennes jusqu'au port d'Anvers) émet un CO\textsubscript{2} par kilogramme de biomasse comparable à celui d'un camion venant d'Allemagne, en raison de l'efficacité énergétique de ces grands navires. Une limite demeure toutefois~: ce calcul ne tient pas compte des autres polluants émis par ces navires, comme les oxydes de soufre.

La quantité de biomasse nécessaire pour alimenter une centrale de taille significative est considérable~: l'ancienne centrale d'Awirs recevait ainsi une à deux péniches par jour.

\subsubsection{Y a-t-il un lien avec le prix de la commode IKEA ?}

Non. La bioénergie ne pèse pas sur les prix des matériaux bois, car elle valorise des \textbf{résidus} de l'industrie forestière qui n'auraient aucun autre usage. Ce ne sont pas des bois d'œuvre susceptibles d'entrer en concurrence avec la production de meubles, de planches ou de papier.

Le professeur Blondeau rappelle que la demande en bioénergie ne tire pas la filière forestière~: c'est la demande en bois à haute valeur ajoutée (construction, papier, carton) qui la structure. De plus, depuis la réduction de la consommation de papier (liée notamment au développement du numérique), les volumes extraits des forêts ont globalement diminué. La part consacrée à la bioénergie reste proportionnellement stable.

Valoriser ces résidus a même un effet positif indirect~: cela apporte un revenu aux gestionnaires forestiers, les incitant à reboiser et à entretenir leurs forêts, plutôt que de les abandonner ou de les convertir à d'autres usages.

\subsubsection{Le problème récent du refus d'importation en Belgique}

Pour bénéficier de subsides (certificats verts), les opérateurs utilisant de la biomasse doivent démontrer que celle-ci est durable, selon des critères définis par l'Union européenne et, le cas échéant, renforcés par les régions. En Flandre, la procédure implique notamment la consultation de fédérations professionnelles du secteur bois (meubles, papier, carton), afin de vérifier que la biomasse importée n'entre pas en concurrence avec leurs approvisionnements.

Dans le cas évoqué (enregistré en janvier 2022, concernant la centrale de Rodenhuize), ces deux fédérations avaient rendu des avis concordants et favorables à l'importation. Malgré cela, la ministre compétente a bloqué le dossier.

Le professeur Blondeau ne se prononce pas sur la décision politique en tant que telle, mais regrette la \textbf{communication} qui en a découlé~: la presse et la ministre elle-même ont évoqué le risque de déforestation en Russie pour brûler des arbres en Belgique, alors que la source concernée avait été certifiée durable par un mécanisme reconnu à l'échelle européenne, et que les fédérations professionnelles potentiellement en concurrence avaient confirmé l'absence de conflit d'usage. Ce cliché, selon lui, ne devait pas être invoqué dans ce contexte précis.

\subsubsection{Les mécanismes de certification du bois}

L'Union européenne définit des critères de durabilité pour la biomasse utilisée à des fins énergétiques, puis reconnaît certains mécanismes de certification comme conformes à ces critères. Parmi les plus connus figurent le \href{https://www.fsc.be/fr-be/faq-01}{FSC (\textit{Forest Stewardship Council})} et le PEFC, initialement développés pour les produits bois (planches, papier, carton), et qui couvrent la gestion forestière durable.

Pour la bioénergie, des critères supplémentaires s'appliquent~:
\begin{itemize}
    \item bilan carbone sur toute la chaîne (collecte, transport, transformation) ;
    \item maintien ou accroissement du stock de carbone forestier ;
    \item respect de la biodiversité ;
    \item légalité des opérations, droits des communautés locales, interdiction du travail des enfants.
\end{itemize}

Des auditeurs indépendants vérifient le respect de ces standards sur le terrain et délivrent des certificats traçables tout au long de la chaîne de valeur. Ces mécanismes impliquent de plus en plus l'ensemble des parties prenantes (\textit{multi-stakeholder})~: producteurs d'énergie, forestiers, négociants, et organisations non gouvernementales.

Le professeur Blondeau souligne que la bioénergie joue un rôle \textbf{pionnier} parmi les énergies renouvelables en matière de certification de durabilité, et que d'autres secteurs (éolien, solaire, batteries) commencent à s'en inspirer. Aucun mécanisme n'est parfait, notamment sur des aspects difficiles à auditer~: impact sur la biodiversité, relations avec les communautés indigènes. Mais le dialogue continu entre acteurs permet d'améliorer progressivement ces standards.

\subsubsection{La biomasse, est-ce durable ?}

Selon les projections de la Commission européenne et de l'\href{https://www.ieabioenergy.com/}{Agence internationale de l'énergie (IEA)}, l'utilisation de la bioénergie devrait \textbf{doubler} d'ici 2050, en cohérence avec les ressources durables estimées disponibles.

Aujourd'hui, la biomasse représente environ \textbf{60~\%} de l'énergie renouvelable consommée en Europe~--- davantage que toutes les autres sources renouvelables réunies (vent, soleil, hydraulique, géothermie). Ce rôle dominant va diminuer en termes relatifs à mesure que l'éolien et le solaire croîtront plus rapidement, mais la biomasse continuera d'augmenter en valeur absolue, pour atteindre une part d'environ 25~\% du mix renouvelable à l'horizon 2050.

La limite fondamentale de la biomasse est la \textbf{disponibilité des ressources durables}~: ce n'est pas la quantité brute de biomasse qui compte, mais la quantité de biomasse certifiée durable.

Sur la question de la neutralité carbone, le professeur Blondeau s'appuie sur un rapport du \href{https://publications.jrc.ec.europa.eu/repository/handle/JRC122719}{Centre commun de recherche de la Commission européenne (JRC)} pour clarifier un débat fréquent. Il existe deux législations distinctes~: l'une portant sur les énergies renouvelables, l'autre sur la foresterie (\href{https://ec.europa.eu/clima/eu-action/forests-and-agriculture/land-use-and-forestry-regulation-2021-2030_en}{LULUCF} --- prononcé \textit{loulousièf} par le Professeur Blondeau). Le LULUCF impose que le stock de carbone forestier soit maintenu ou augmenté~: dès qu'un arbre est extrait, ce carbone est comptabilisé comme émis. Dès lors, comptabiliser à nouveau le CO\textsubscript{2} en sortie de cheminée reviendrait à le compter deux fois. La biomasse est donc considérée comme neutre en CO\textsubscript{2} non pas par convention politique arbitraire, mais parce que le carbone est déjà comptabilisé au niveau de la forêt.

Le JRC illustre cela par une analogie~: plutôt que de suivre les enfants dans chaque magasin, les parents vérifient le solde du compte bancaire. Si ce solde est stable ou croît, les dépenses sont maîtrisées. Ce découplage implique cependant une bonne communication entre les gestionnaires forestiers et les autorités, pour éviter qu'une extraction excessive ne soit détectée trop tardivement.

Enfin, si la biomasse favorise le reboisement (en rendant la filière forestière plus rentable) et si on y ajoute une capture du CO\textsubscript{2} en sortie de centrale (\textit{BECCS}), les émissions nettes peuvent devenir négatives~--- un atout considérable pour atteindre la neutralité climatique.

\subsubsection{Le trilemme de la biomasse}

Le trilemme énergétique désigne l'impossibilité, pour toute source d'énergie actuelle, de satisfaire simultanément trois critères~:
\begin{enumerate}
    \item \textbf{Disponibilité}~: géopolitique (indépendance vis-à-vis d'un pays ou d'une région) et temporelle (production à la demande) ;
    \item \textbf{Impact environnemental}~: faibles émissions de CO\textsubscript{2} et de polluants locaux ;
    \item \textbf{Coût}~: accessibilité économique sans subsides excessifs.
\end{enumerate}

La biomasse se positionne favorablement sur ces trois axes~:
\begin{itemize}
    \item \textbf{Disponibilité}~: présente sur plusieurs continents, elle réduit la dépendance énergétique. Elle est aussi disponible à la demande, car stockable, contrairement au vent ou au soleil.
    \item \textbf{Environnement}~: neutre en CO\textsubscript{2} si certifiée durable. Les émissions de particules et d'oxydes d'azote restent un point d'attention, surtout aux petites échelles.
    \item \textbf{Coût}~: comparable aux sources fossiles traditionnelles. Des subsides restent souvent nécessaires, en partie parce que les installations existantes sont parfois vétustes et peu efficaces. Des centrales modernes dédiées à la biomasse amélioreraient sensiblement ce bilan.
\end{itemize}

La biomasse présente également un avantage souvent sous-estimé~: elle génère de l'emploi local, depuis la collecte jusqu'à la conversion, ce qui la distingue favorablement de certaines technologies renouvelables dont la chaîne de valeur est largement externalisée.

\subsubsection{Les sources de chaleur verte}

L'être humain utilise la biomasse comme source de chaleur depuis près d'un million d'années. Aujourd'hui, elle reste la principale source de chaleur renouvelable, mais d'autres options existent.

L'IEA distingue la \textbf{biomasse traditionnelle} (foyers ouverts, faible rendement, émissions importantes) de la \textbf{biomasse moderne} (installations performantes, rendement 10 à 100 fois supérieur, émissions réduites). La transition de l'une vers l'autre constitue un levier majeur de décarbonation, notamment pour le chauffage domestique en ville.

Les autres sources de chaleur verte incluent~:
\begin{itemize}
    \item les \textbf{pompes à chaleur}~: très efficaces pour les bâtiments bien isolés avec chauffage basse température (30~°C ou moins), mais inadaptées aux hautes températures industrielles (100--400~°C) ;
    \item les \textbf{panneaux solaires thermiques}~: source décentralisée, complémentaire ;
    \item l'\textbf{hydrogène} et les \textbf{électrofuels}~: vecteurs de chaleur produits à partir d'électricité renouvelable, encore en développement.
\end{itemize}

La biomasse présente un avantage distinctif~: elle peut fournir de la chaleur à \textbf{haute température}, y compris pour des usages industriels, là où les pompes à chaleur ne sont pas compétitives. Elle est aussi compatible avec les systèmes de chauffage existants, ce qui facilite une transition progressive sans rénovation lourde du bâti.

\subsubsection{Les réseaux de chaleur, intéressant en Belgique ?}

Un réseau de chaleur consiste à distribuer de la chaleur sous forme d'eau chaude ou de vapeur via des canalisations enterrées, depuis une ou plusieurs sources de production vers les utilisateurs (bâtiments résidentiels, bureaux, industries, équipements publics). La source peut être la biomasse, un incinérateur de déchets, des panneaux solaires thermiques, une pompe à chaleur, ou une combinaison de ces options. L'utilisateur final est indépendant de la source~: seul le tuyau d'eau chaude change.

Cela confère aux réseaux de chaleur une grande \textbf{flexibilité}~: si une meilleure technologie de production émerge dans 20 ans, seule la centrale doit être remplacée, sans toucher aux installations chez chaque abonné.

Les pays scandinaves sont très avancés dans ce domaine. En Belgique, le développement reste limité, pour plusieurs raisons~:
\begin{itemize}
    \item un réseau n'est économiquement pertinent que si la \textbf{densité de demande} est suffisante et la \textbf{distance} entre utilisateurs faible (sinon, les pertes en ligne et le surdimensionnement de la centrale pèsent trop sur le bilan) ;
    \item les promoteurs privés préfèrent souvent des solutions plus simples à vendre (pompe à chaleur individuelle par logement) plutôt que d'assumer la conception, le financement et l'exploitation d'un réseau sur 25 ans ;
    \item les autorités publiques manquent d'outils d'aide à la décision pour évaluer cas par cas la pertinence d'un réseau de chaleur.
\end{itemize}

Les cas les plus favorables sont les \textbf{nouveaux quartiers denses}, où les travaux de voirie sont déjà programmés et où la mixité des usages (logements, bureaux, commerces, équipements publics) lisse la demande sur la journée, améliorant la rentabilité de l'infrastructure. Le groupe de recherche du professeur Blondeau travaille au développement d'outils permettant aux décideurs publics d'évaluer ces projets de manière éclairée, en estimant le coût de la chaleur et le soutien public nécessaire au cas par cas.

\subsubsection{Le stockage d'énergie, important dans un monde renouvelable}

Dans un système dominé par des sources intermittentes (vent et soleil), trois leviers permettent d'assurer l'équilibre entre offre et demande~:
\begin{enumerate}
    \item \textbf{La gestion de la demande}~: décaler la consommation (machine à laver, recharge de véhicule) aux moments où l'électricité est disponible et bon marché ;
    \item \textbf{Les centrales de backup}~: centrales thermiques (gaz, biomasse) démarrables à la demande pour couvrir les périodes sans vent ni soleil (c'est l'objet du CRM belge) ;
    \item \textbf{Le stockage}~: produire de l'énergie quand elle est disponible et la restituer plus tard.
\end{enumerate}

Le stockage électrique par batteries, bien que prometteur, reste coûteux et limité en capacité. Une alternative est le \textbf{stockage chimique}~: utiliser l'électricité renouvelable excédentaire pour produire des molécules énergétiques via électrolyse. La première brique est l'\textbf{hydrogène} (H\textsubscript{2}), obtenu par électrolyse de l'eau. Il peut ensuite être combiné à du carbone (fourni notamment par la biomasse, source de carbone renouvelable) pour produire des \textbf{électrofuels}~: méthane de synthèse, ammoniac, éthanol, etc. Ces vecteurs sont plus faciles à stocker et à transporter que l'hydrogène pur.

Le rendement de la chaîne électrolyse-stockage-reconversion n'est pas optimal, mais le professeur Blondeau souligne que la \textbf{valeur économique de l'énergie} ne dépend pas uniquement de son rendement de production~: dans un monde 100~\% renouvelable, l'énergie disponible au bon moment aura une valeur élevée, justifiant ces conversions intermédiaires. De plus, l'hydrogène produit à partir de grandes installations hydroélectriques présente un bilan énergétique global plus favorable.

La biomasse, stockable par nature, joue ici un rôle de \textbf{complément indispensable} au vent et au soleil, en tant que source d'énergie pilotable et de carbone renouvelable pour la production d'électrofuels.